\chapter{The Vol Surface and Vol Trading}\label{ch:iv2:vol-surface-and-vol-trading}

\begin{ivtrains}
Reading a surface, moving it the right way when the stock moves, and trading volatility with variance, dispersion and events. Unless a stem says otherwise, rates are zero, and 252 trading days make a year. All levels here are invented for the arithmetic.
\end{ivtrains}

\section{Reading a surface}

\begin{question}[$\star$ \iqroles{QT, Trader, QR, Risk} \iqfirm{any} \iqlevel{junior}]\label{iq:iv2:vol-surface-and-vol-trading:1}
A stock's five daily returns are $+1\%$, $-1\%$, $+2\%$, $-2\%$ and $0\%$. What is its annualised realised volatility, taking the mean return as zero? Why do desks usually take the mean as zero for daily data?
\end{question}

\begin{question}[$\star$ \iqroles{QT, Trader, QR, Risk} \iqfirm{any} \iqlevel{junior}]\label{iq:iv2:vol-surface-and-vol-trading:2}
A stock has 24\% implied volatility. Roughly how much is it expected to move in a day, by the rule of 16? Why 16? What would the same figure be if you count every calendar day?
\end{question}

\begin{question}[$\star$ \iqroles{QT, Trader, QR} \iqfirm{any} \iqlevel{junior}]\label{iq:iv2:vol-surface-and-vol-trading:3}
You are long an at-the-money option, delta-hedged, on a \$100 stock with 25\% implied volatility, and you hold it one trading day. How large a move do you need to break even? How often will the stock beat that, if the model holds?
\end{question}

\begin{question}[$\star$ \iqroles{QT, Trader, QR, Risk} \iqfirm{any} \iqlevel{junior}]\label{iq:iv2:vol-surface-and-vol-trading:4}
At-the-money implied volatility is 25\% for three months and 22\% for six months. What is the forward volatility from three to six months? Why do you average variances and not volatilities?
\end{question}

\begin{question}[$\star$ \iqroles{QT, Trader, QR} \iqfirm{market maker, hedge fund, bank} \iqlevel{junior}]\label{iq:iv2:vol-surface-and-vol-trading:5}
A \$100 stock has at-the-money implied volatility of 30\% at three months and 20\% at six months. What is wrong? What do you trade, and what does it pay you up front?
\end{question}

\begin{question}[$\star$ \iqroles{QT, Trader, QR, Analyst} \iqfirm{any} \iqlevel{junior}]\label{iq:iv2:vol-surface-and-vol-trading:6}
On an equity index, a one-year put 10\% below the stock has an implied volatility of 26\%, the at-the-money option has 20\%, and the call 10\% above has 17\%. Why does the surface look like that when the Black-Scholes model says it should be flat? Give your best three reasons, and say which one you would lead with.
\end{question}

\begin{question}[$\star$ \iqroles{QT, Trader, QR} \iqfirm{market maker, hedge fund, bank} \iqlevel{junior}]\label{iq:iv2:vol-surface-and-vol-trading:7}
A variance swap has a vega notional of \$100\,000 and a strike of 20. What is the variance notional? What do you make if you are long and realised volatility is 25? What if it is 15? Why is the gain larger than the loss for the same five points?
\end{question}

\begin{question}[$\star$ \iqroles{QR, QT, Risk} \iqfirm{any} \iqlevel{junior}]\label{iq:iv2:vol-surface-and-vol-trading:8}
You measure a stock's realised volatility as 20\% from 20 daily returns. Roughly what is the standard error of that number if returns are normal? How many days do you need to bring it down to 1 point?
\end{question}

\begin{question}[$\star$ \iqroles{QT, Trader, QR} \iqfirm{hedge fund, market maker, bank} \iqlevel{junior}]\label{iq:iv2:vol-surface-and-vol-trading:9}
An index with many equally weighted members has an implied volatility of 18\%, and each member has one of 30\%. What is the index's implied correlation? What does a rise in that number tell you?
\end{question}

\section{Smiles and how they move}

\begin{question}[$\star\star$ \iqroles{QR, QT} \iqfirm{hedge fund, market maker, bank} \iqlevel{junior, experienced}]\label{iq:iv2:vol-surface-and-vol-trading:10}
You believe a \$100 stock's volatility for the next year is either 10\% or 30\%, with equal chances, and you will know it right away. Price calls as the average of the two Black-Scholes prices. What implied volatility does that give at strikes of 100, 80 and 130? Why is the at-the-money figure near the average of the two, and not their root mean square?
\end{question}

\begin{question}[$\star\star$ \iqroles{QR, QT, QD} \iqfirm{any} \iqlevel{junior, experienced}]\label{iq:iv2:vol-surface-and-vol-trading:11}
For a \$100 stock with a one-year expiry and a zero rate, you see calls at 95, 100 and 105 worth 10.52, 7.97 and 5.91 when priced at 20\% volatility. Give the market's probability density at 100 from these three prices. What is the true lognormal density there? What would a negative answer mean?
\end{question}

\begin{question}[$\star\star$ \iqroles{QT, Trader, QR} \iqfirm{market maker, hedge fund, bank} \iqlevel{junior, experienced}]\label{iq:iv2:vol-surface-and-vol-trading:12}
At a stock price of 100, the at-the-money one-year implied volatility is 20\%, and the smile falls one point for each 10\% of strike. The stock jumps to 105. What is the at-the-money volatility, and what is the 100 strike's volatility, under sticky strike and under sticky moneyness? You are long the 100 call, with a vega of 0.39 per point: what is the difference in your profit?
\end{question}

\begin{question}[$\star\star$ \iqroles{QT, Trader, QR, Risk} \iqfirm{market maker, hedge fund, bank} \iqlevel{junior, experienced}]\label{iq:iv2:vol-surface-and-vol-trading:13}
An at-the-money one-year call on a \$100 stock has a Black-Scholes delta of 0.54 and a vega of 39.7 per unit of volatility, and the smile's slope is $-0.002$ per dollar of strike. What delta do you use under sticky strike, under sticky moneyness and under a local volatility model that moves the smile the other way? Which one is right?
\end{question}

\begin{question}[$\star\star$ \iqroles{QT, Trader} \iqfirm{market maker, hedge fund, bank} \iqlevel{junior, experienced}]\label{iq:iv2:vol-surface-and-vol-trading:14}
The one-month smile has a slope of $-0.5$ points of volatility per 1\% of strike. A stock falls 2\% and at-the-money implied volatility rises 1.5 points. How much of that is the stock sliding along the smile, and how much is the smile moving? What did the option struck at the old at-the-money level do?
\end{question}

\begin{question}[$\star\star$ \iqroles{QT, Trader, QR} \iqfirm{hedge fund, market maker} \iqlevel{junior, experienced}]\label{iq:iv2:vol-surface-and-vol-trading:15}
You sell a one-year variance swap on an index at a strike of 22, with a vega notional of \$100\,000, and it realises 18. What do you make? Roughly how many points of vega is that? Why might implied stay above realised on average, and what is the catch in selling it?
\end{question}

\begin{question}[$\star\star$ \iqroles{QR, QT} \iqfirm{hedge fund, market maker} \iqlevel{junior, experienced}]\label{iq:iv2:vol-surface-and-vol-trading:16}
You have open, high, low and close data. How do you estimate a day's variance from the range? Where does the factor $1/(4\ln 2)$ come from, and how much better is it than a squared close-to-close return? What biases it in real data?
\end{question}

\begin{question}[$\star\star$ \iqroles{QT, Trader, QR} \iqfirm{market maker, hedge fund} \iqlevel{junior, experienced}]\label{iq:iv2:vol-surface-and-vol-trading:17}
It is Friday's close. A stock has a business-time volatility of 25\%, counted on 252 trading days. An option expires at Monday's close, which is one trading day and three calendar days away. What volatility does a screen show, using 365 calendar days? Why does it look so low? What does this do to a trader looking at the term structure?
\end{question}

\begin{question}[$\star\star$ \iqroles{QT, Trader, QR, Analyst} \iqfirm{market maker, hedge fund} \iqlevel{junior, experienced}]\label{iq:iv2:vol-surface-and-vol-trading:18}
Earnings fall within the next five trading days. The five-day option has an implied volatility of 60\%, and the 25-day option 40\%. Neither has another event in it. What are the base volatility and the implied move on the day, as one standard deviation? What is the expected absolute move? What do you assume?
\end{question}

\begin{question}[$\star\star$ \iqroles{QT, Trader, Analyst} \iqfirm{market maker, hedge fund} \iqlevel{junior, experienced}]\label{iq:iv2:vol-surface-and-vol-trading:19}
A \$100 stock has a base volatility of 25\% and an earnings move of one standard deviation of 6\% tomorrow, and 30 calendar days to expiry, on a 365-day year. What is the implied volatility of the at-the-money option? You buy the straddle at about \$7.46. After the event, implied volatility is 25\%, and 29 days remain. What is the straddle worth if the stock is unchanged, or at 103, 106 or 110? How far must it move for you to break even?
\end{question}

\begin{question}[$\star\star$ \iqroles{QT, QR} \iqfirm{hedge fund, market maker} \iqlevel{junior, experienced}]\label{iq:iv2:vol-surface-and-vol-trading:20}
Next year's realised volatility has a mean of 20\% and a standard deviation of 5\%. What is the fair variance swap strike? What is the fair volatility swap strike? Why does a volatility swap pay you less in a big move? Why is it hard to hedge?
\end{question}

\begin{question}[$\star\star$ \iqroles{QT, QR, Trader} \iqfirm{hedge fund, market maker, bank} \iqlevel{junior, experienced}]\label{iq:iv2:vol-surface-and-vol-trading:21}
An index has 50 equally weighted stocks, each with 30\% volatility and a pairwise correlation of 0.3. What is the index's volatility? You sell index variance and buy variance on all the stocks. If correlation realises at 0.2 and stock volatility at 30\%, how does that help you? What if stock volatility is 33\%? What is this trade called, and what risk are you really holding?
\end{question}

\begin{question}[$\star\star$ \iqroles{QT, Trader, QR} \iqfirm{market maker, hedge fund, bank} \iqlevel{junior}]\label{iq:iv2:vol-surface-and-vol-trading:22}
At a one-year expiry on a \$100 stock, the at-the-money volatility is 20\% and the 90 strike's is 25\%. What does a 100-90 put spread cost? What would it cost if you priced both at 20\%? Why is the spread cheaper with skew? Who is therefore paying what?
\end{question}

\begin{question}[$\star\star$ \iqroles{QR, QD} \iqfirm{hedge fund, bank, market maker} \iqlevel{junior, experienced}]\label{iq:iv2:vol-surface-and-vol-trading:23}
A Heston model has a variance reversion speed of 2, a long-run variance of 0.04 and a vol of vol of 0.5. Does it meet the Feller condition? What does it mean if it does not, and what do you do about it when you simulate it?
\end{question}

\begin{question}[$\star\star$ \iqroles{QR, QD, Risk} \iqfirm{any} \iqlevel{junior, experienced}]\label{iq:iv2:vol-surface-and-vol-trading:24}
You need the two-month at-the-money volatility from one-month at 20\% and three-month at 22\%. How do you interpolate? What do you get? Why not interpolate linearly in volatility? How do you make sure that no calendar arbitrage comes in?
\end{question}

\begin{question}[$\star\star$ \iqroles{QT, QR, Risk} \iqfirm{hedge fund, market maker, bank} \iqlevel{junior, experienced}]\label{iq:iv2:vol-surface-and-vol-trading:25}
You are long a one-year variance swap with a vega notional of \$100\,000 and a strike of 20. After six months, the realised volatility so far is 25, and the implied volatility for the rest is 22. What is it worth? Why is this mark so simple? How does its vega compare with what it was at the start?
\end{question}

\section{Variance and models}

\begin{question}[$\star\star\star$ \iqroles{QR, QT} \iqfirm{hedge fund, market maker} \iqlevel{junior, experienced}]\label{iq:iv2:vol-surface-and-vol-trading:26}
Why does a variance swap's replication use options weighted by one over the strike squared? Show, for a stock at 80, 100 and 120 and 20\% flat volatility over one year, that this strip has the same dollar gamma at each, and that a strip with equal weights does not. What is the fair strike, and what does this have to do with the log contract?
\end{question}

\begin{question}[$\star\star\star$ \iqroles{QR} \iqfirm{hedge fund, bank} \iqlevel{experienced}]\label{iq:iv2:vol-surface-and-vol-trading:27}
A vendor's one-year smile rises steeply in the far wings, and it has a total variance of $3|\ln(K/S)|$ for large strikes. Show from the call price that this is impossible. What is the largest slope allowed? For a stock at 100, what is the largest implied volatility at a strike of 400, if the bound applies?
\end{question}

\begin{question}[$\star\star\star$ \iqroles{QR, QD} \iqfirm{hedge fund, bank, market maker} \iqlevel{junior, experienced}]\label{iq:iv2:vol-surface-and-vol-trading:28}
Write down Dupire's formula for local volatility with a zero rate. Check it for a flat 20\% surface. Then check it for the mixture where calls are the average of Black-Scholes prices at 10\% and at 30\%, at a one-year strike of 100. What do you get? Why is it so different from the implied 19.98\%, and how would you read it?
\end{question}

\begin{question}[$\star\star\star$ \iqroles{QR, QT} \iqfirm{hedge fund, market maker, bank} \iqlevel{junior, experienced}]\label{iq:iv2:vol-surface-and-vol-trading:29}
A \$100 stock has a local volatility of $20\% - 0.2\ln(S/100)$, with a zero rate. What is the implied volatility at one year for strikes of 80, 100 and 120? Compare the implied slope in log strike with the local one. What is the rule, and why does it work?
\end{question}

\begin{question}[$\star\star\star$ \iqroles{QR, QT} \iqfirm{hedge fund, market maker, bank} \iqlevel{experienced}]\label{iq:iv2:vol-surface-and-vol-trading:30}
A stock has a local volatility of $20\% - 0.2\ln(S/100)$, with a zero rate. With the stock at 100, the one-year implied volatility at a strike of 100 is 19.98\%. With the stock now at 105, it is 19.01\% at a strike of 105 and 19.48\% at 100. What did the at-the-money volatility do, and which way did the smile move? How does that compare with what markets tend to do? What does it do to your hedge?
\end{question}

\begin{question}[$\star\star\star$ \iqroles{QR, QD} \iqfirm{hedge fund, bank, market maker} \iqlevel{junior, experienced}]\label{iq:iv2:vol-surface-and-vol-trading:31}
In Heston's model, the variance now is 0.04 and its long-run level 0.09, with a reversion speed of 2. What does the at-the-money term structure look like at three months, one year and four years, roughly, before any effect of vol of vol? Now set the variance and the long-run level both to 0.04, keep the speed of 2, and take a vol of vol of 0.5, on a \$100 stock with a zero rate. The one-year implied volatility, in percent, at strikes of 80, 100 and 120 is 23.1, 18.1 and 14.6 at a correlation of $-0.7$, and 20.4, 18.8 and 19.9 at a correlation of zero. What sets the slope and what the curvature? Why is the at-the-money below 20\%?
\end{question}

\begin{question}[$\star\star\star$ \iqroles{QR, QT} \iqfirm{bank, hedge fund} \iqlevel{junior, experienced}]\label{iq:iv2:vol-surface-and-vol-trading:32}
In the lognormal SABR model, with an initial volatility of 20\%, a vol of vol of 0.4, a correlation of $-0.3$ and one year, what is the at-the-money implied volatility by Hagan's formula? What does it give at 80 and 120? How would you check it? What is SABR's advantage over local volatility, and its well-known weakness?
\end{question}

\begin{question}[$\star\star\star$ \iqroles{QR, QT} \iqfirm{hedge fund, market maker, bank} \iqlevel{junior, experienced}]\label{iq:iv2:vol-surface-and-vol-trading:33}
A \$100 stock, with a zero rate, has a diffusion volatility of 15\%, and it can gap down 20\%, at a rate of 0.1 times a year. Write its call price as a Poisson mixture of Black-Scholes prices, with the right drift. What are the implied volatilities at 90, 100 and 110, at one month and at one year? Why is the short-dated skew so steep? What can a pure diffusion model not do here?
\end{question}

\begin{question}[$\star\star\star$ \iqroles{QR, QT} \iqfirm{hedge fund, market maker} \iqlevel{experienced}]\label{iq:iv2:vol-surface-and-vol-trading:34}
You have sold a variance swap and hedged it with the log-contract strip. The stock jumps down 20\% in one step, or up 20\%. What does the swap's squared log return pay in each case, and what does your strip hedge give? Who loses, and how much per unit of variance? What does this imply about the strike?
\end{question}

\begin{question}[$\star\star\star$ \iqroles{QR, QT} \iqfirm{hedge fund, market maker} \iqlevel{junior, experienced}]\label{iq:iv2:vol-surface-and-vol-trading:35}
A one-year smile is $\sigma(K) = 20\% - 0.1\ln(K/100)$, with a floor of 3\%, on a \$100 stock with a zero rate. By the strip of out-of-the-money options, what is the fair variance swap strike? A desk's rule is $\sigma_{\text{ATM}}^2(1 + 3T\,b^2)$, with slope $b$ per unit of log strike. What does the rule give? Why is the strike above at-the-money?
\end{question}

\begin{question}[$\star\star\star$ \iqroles{QT, QR} \iqfirm{hedge fund, market maker, bank} \iqlevel{junior, experienced}]\label{iq:iv2:vol-surface-and-vol-trading:36}
Stocks have 30\% volatility. An index of very many of them is at 100, with zero rates. Its correlation for the year is 0.2 with probability 4/5 and 0.8 with probability 1/5. What is the index's volatility in each state? What is its one-year smile at strikes of 80, 100 and 120, if the two states have the same forward? Does that give a skew? What feature must you add to get one, and why are index skews steeper than single-stock ones?
\end{question}

\begin{question}[$\star\star\star$ \iqroles{QT, Trader, QR} \iqfirm{hedge fund, market maker} \iqlevel{junior, experienced}]\label{iq:iv2:vol-surface-and-vol-trading:37}
You sell an at-the-money one-month straddle on a \$100 stock at 20\%, and the stock's volatility really is 16\%. You hedge daily, 21 times. What is the mean and the standard deviation of your profit? How often do you lose, and what is the 1st percentile? What does this have to say about the size to trade, when real distributions have jumps?
\end{question}

\section{Three to build}

\begin{question}[$\star\star\star\star$ \iqroles{QR, QD} \iqfirm{hedge fund, market maker} \iqlevel{experienced}]\label{iq:iv2:vol-surface-and-vol-trading:38}
Derive the variance swap's fair strike from the log contract, and write it as a strip of options that are out of the money. Compute it for the smile $\sigma(K) = 20\% - 0.1\ln(K/100)$, with a floor of 3\%, on a \$100 stock with a zero rate over a year, using strikes from 50 to 200 in steps of \$1, \$5 and \$10, and from 70 to 140 in steps of \$5. Compare with the exact integral. Which error is which? Which matters more in a stress?
\end{question}

\begin{question}[$\star\star\star\star$ \iqroles{QR, QD} \iqfirm{hedge fund, bank} \iqlevel{experienced}]\label{iq:iv2:vol-surface-and-vol-trading:39}
Show that, if volatility is random and independent of the stock's own noise, a call's price is the average of Black-Scholes prices over the distribution of realised variance. Test it on a Heston model for a \$100 stock over one year with a zero rate, $v_0 = \theta = 0.04$, $\kappa = 2$, a vol of vol of 0.5 and zero correlation, against a plain simulation of the same model at strikes of 80, 100 and 120. How much does the standard error fall? What happens to the argument with correlation?
\end{question}

\begin{question}[$\star\star\star\star$ \iqroles{QR, QD} \iqfirm{hedge fund, bank} \iqlevel{experienced}]\label{iq:iv2:vol-surface-and-vol-trading:40}
A \$100 stock with a zero rate has, at every strike and maturity, a call price equal to the equal-weight average of the Black-Scholes prices at 10\% and at 30\% volatility. Find the local volatility surface in closed form as an average of the two variances, with the weights given by each state's density. Reprice one-year calls at 70, 100 and 130 in that local volatility model. Does a plain daily Euler simulation get them? What goes wrong near time zero, and how do you solve it properly?
\end{question}
